% !PRINTED 357
\weberpage{357}
%<<CONT
are relatively prime, $x$ cannot be divisible by $p$, and $\mathfrak{a}$
must be divisible by $p$.

If we put $\mathfrak{a} = p^k \mathfrak{a}_1$, where $\mathfrak{a}_1$ is no longer divisible by $p$,
and denote by $\mathfrak{p}$ a prime ideal dividing $p$, then by
(11) $\mathfrak{p}$ must divide $\mathfrak{m}$ or $\mathfrak{m}'$, and since $\mathfrak{m}$ was assumed to be primary, $\mathfrak{p}$ cannot be $= p$. Therefore $N(\mathfrak{p}) = p = \mathfrak{p}\mathfrak{p}'$,
and if $\mathfrak{m}$ is divisible by $\mathfrak{p}$, it cannot at the same time be divisible by $\mathfrak{p}'$.
Consequently $\mathfrak{m}$ is divisible by $\mathfrak{p}^k$. Hence, by
the definition of $\mathfrak{m}$ (Theorem 6), $\dfrac{b+\sqrt D}{2}$ is also divisible by $\mathfrak{p}^k$,
and consequently $\dfrac{b+\sqrt D}{2} y$ at least by $\mathfrak{p}^{k+1}$, while $ax$ and consequently $\mu$
exactly by $\mathfrak{p}^k$, whereas on the other hand $\mathfrak{a}\mathfrak{m} = \mu$ is likewise divisible at least by $\mathfrak{p}^{k+1}$, in which lies a contradiction. The like holds also when $p = \mathfrak{p}^2$, that is, in $\Delta$ (but not in $Q$).
Consequently no rational prime can be contained in $\mathfrak{a}$,
and $\mathfrak{a}$ is primary.

From what has been proved so far we can formulate the following
theorem:

9. \gesperrt{Theorem. If}
\[
\mu = ax + \frac{b+\sqrt D}{2} y = \mathfrak{m}\mathfrak{a},
\]
\[
\mu' = ax' + \frac{b+\sqrt D}{2} y' = \mathfrak{m}\mathfrak{a}'
\]
% Source scan supports b before \sqrt D, matching the parallel line above.
\gesperrt{are two numbers coprime to $Q$, if neither $x$ and $y$
nor $x'$ and $y'$ have a common divisor, if}
\[
N(\mathfrak{a}) = N(\mathfrak{a}') = A,
\]
\gesperrt{and if for one and the same number $B$}
\[
\begin{aligned}
\frac{B+\sqrt D}{2}&\equiv 0 \pmod{\mathfrak{a}}\\
&\equiv 0 \pmod{\mathfrak{a}'},
\end{aligned}
\]

\begin{center}\rule{0.2\linewidth}{0.4pt}\end{center}
\gesperrt{then the two ideals $\mathfrak{a}$ and $\mathfrak{a}'$ are identical, and the
numbers $\mu$, $\mu'$ differ only by a unit factor.}

For $\mathfrak{a}$ as well as $\mathfrak{a}'$ are primary by 8., and by 6. are
defined as the greatest common divisor of
\[
A \text{ and } \frac{B+\sqrt D}{2}.
\]
